\documentclass[10pt,a4paper]{amsart}
\usepackage[margin=19mm]{geometry}
\usepackage{amsmath,amssymb,mathtools}
\usepackage[scr=rsfs,cal=euler]{mathalfa}
\usepackage{xfrac}
\usepackage[unicode,hidelinks,bookmarks=false]{hyperref}
\newcommand{\ie}{i.e.\ }
\newcommand{\cf}{cf.\ }
\newcommand{\resp}{respectively\ }
\newcommand{\Subsection}{\subsection}
\providecommand{\nobreakdash}{\nobreak\hbox{-}}
\setlength{\emergencystretch}{2em}
\multlinegap0pt
\usepackage{amssymb}
\newtheorem{proposition}{Proposition}
\newcommand{\Z}{\mathbb{Z}}
\title{On semigroups that are prime in the sense of Tarski, and groups prime in the senses of Tarski and of Rhodes}
\author{George M. Bergman}
\date{Source: arXiv:2409.15541v3 (CC BY 4.0)}
\begin{document}
\maketitle

\noindent\emph{Source and licence.} On semigroups that are prime in the sense of Tarski, and groups prime in the senses of Tarski and of Rhodes. Version arXiv:2409.15541v3 (CC BY 4.0), distributed by arXiv under the Creative Commons Attribution 4.0 International licence, http://creativecommons.org/licenses/by/4.0/, as recorded in the arXiv licence metadata for this exact version and shown on the abstract page https://arxiv.org/abs/2409.15541v3. Full licence notice: \texttt{licenses/arxiv-2409.15541-CC-BY-4.0.txt}.\par\smallskip

\noindent\textit{Editorial scope.} Counted unit: the article's proposition P.x-prime (source0.tex lines 1596--1609). The article's complete original proof of it is delivered with it (source0.tex lines 1611--1734, one \texttt{proof} environment of 124 lines). No mathematics is written, repaired or re-derived: the statement and the proof are the article's own lines, in the article's own notation, and no retained line is substituted or re-flowed.  No further source-local context is needed: every cross-reference of every retained line resolves inside the two delivered spans.  Nothing was omitted from any retained span, so the package carries no omission record.  No retained line contains a citation, so the wrapper carries no bibliography and no bibliography entry.  No retained line contains a figure environment, an image-inclusion command or drawing code, so no image is delivered and the package declares no asset dependency.  Editorial formatting only, in addition: an AMS \texttt{amsart} preamble that restates the article's own declarations and macros used by the retained lines, namely the environments \texttt{proposition} on separate counters, together with the standard packages those lines need.  The retained environments print in this document as Proposition 1 in order of appearance, whereas the article prints the same items with the numbering of its own full document.  The complete native source of the article as distributed in the arXiv e-print for this exact version is bundled unchanged as \texttt{sources/165709/source0.tex}.
\begin{proposition}\label{P.x-prime}
Let $G$ be a monolithic finite group, with monolith~$M.$

Then $G$ is Rhodes-prime with respect to direct products
in the category of finite groups if any
of the following conditions holds:

\hspace{.5em}{\rm(i)} $M$ is noncommutative, or

\hspace{.25em}{\rm(ii)} $G$ is a semidirect product
$M\rtimes K$ of $M$ with a subgroup $K<G,$ or

{\rm(iii)} $G$ is a cyclic group of prime-power order, $\Z_{p^n}.$
\end{proposition}

\begin{proof}
Given a finite monolithic
group $G$ that is {\em not} Rhodes-prime with respect to
direct products, we shall show that neither~(i) nor~(ii) can hold.
The proof that~(iii) implies Rhodes-primeness will be more
straightforward.

To get the first two results, suppose $G$
is not Rhodes-prime with respect to direct products, i.e.,
can be written as a homomorphic image $f(H)$ of a subgroup $H$ of
a direct product $G_0\times G_1$ of finite groups, such that
$G$ is not a subquotient of $G_0$ or of $G_1.$
We can clearly replace $G_0$ and $G_1$ by
the subgroups given by the projections of $H$ onto those two factors,
so assume that those projections are
surjective; i.e., that $H$ is a subdirect product of $G_0$ and $G_1.$
Let us further assume that, for the given
group $G,$ the groups $G_0,$ $G_1,$ and $H$ are chosen so
as to minimize the order of $H.$

Now let
%
\begin{equation}\begin{minipage}[c]{35pc}\label{d.H12}
$H_0 = \{h\in G_0\,|\,(h,e)\in H\}$\hspace{.8em}and\hspace{.8em}\
$H_1 = \{h\in G_1\,|\,(e,h)\in H\}.$
\end{minipage}\end{equation}

If $H_0$ were trivial, then
no two elements of $H$ would fall together under the
projection $G_0\times G_1\to G_1,$ so $G_1$ would be
isomorphic to $H,$ so as $G$ is a homomorphic image
of $H,$ this would contradict the assumption that $G$
was not a subquotient of $G_1.$
So $H_0$ is nontrivial; and similarly $H_1.$
From our assumption that $H$ projects surjectively to each of
$G_0$ and $G_1$ it is also easy to see
that $H_0$ is normal in $G_0,$ and $H_1$ in $G_1.$
(E.g., given $h\in H_0$ and $g_0\in G_0,$ we can find $g_1\in G_1$
such that $(g_0,g_1)\in H;$ and conjugating $(h,e)$ by $(g_0,g_1),$ we
effectively conjugate $h$ by $g_0.)$
Now if $\mathrm{ker}(f)$ had
nontrivial intersection with the subgroup $H_0\times\{e\}$ of $H,$ then
since $H_0\times\{e\}$ and $\mathrm{ker}(f)$ are both normal in $H,$
their intersection would be normal, and dividing $G_0$ by the
projection of this intersection, we could decrease the order of $H.$
This, and the corresponding observation for $\{e\}\times H_1,$ give
%
\begin{equation}\begin{minipage}[c]{35pc}\label{d.triv_cap}
$\mathrm{ker}(f)$ has trivial intersections with $H_0\times\{e\}$
and with $\{e\}\times H_1.$
\end{minipage}\end{equation}

Now let $M_0\times\{e\}$ be a nontrivial subgroup of
$H_0\times\{e\}$ minimal for being normal in $H.$
Since $M_0\subseteq H_0,$~\eqref{d.triv_cap} tells us
that $M_0\times\{e\}$ maps one-to-one into $G,$ hence its
image is a minimal normal subgroup of $G;$ hence that image is~$M.$
Combining with the corresponding observation about
a minimal normal subgroup $\{e\}\times M_1$ of
$\{e\}\times H_1,$ we get
%
\begin{equation}\begin{minipage}[c]{35pc}\label{d.M1,M2}
$H_0\times \{e\}$ is monolithic with monolith $M_0\times\{e\},$
and $\{e\}\times H_1$ is monolithic with monolith
$\{e\}\times M_1,$ and both of these map
isomorphically to $M$ under $f.$
\end{minipage}\end{equation}

But $M_0\times\{e\}$ and $\{e\}\times M_1$
centralize one another in $H,$ so $M$ must be
self-centralizing in $G,$ i.e., commutative,
giving the desired contradiction to~(i).

To get the next assertion, suppose as in~(ii)
that $G$ is a semidirect product $M\rtimes K$ for some $K<G.$
Since $K$ has trivial intersection with the monolith $M$ of $G,$
%
\begin{equation}\begin{minipage}[c]{35pc}\label{d.K_has_no_nm}
$K$ contains no nontrivial normal subgroup of~$G.$
\end{minipage}\end{equation}
%
On the other hand, since $M$ is invariant under conjugation by
members of $K,$ the centralizer of $M$ in $K$ must also be invariant
under conjugation by members of $K;$ and by definition that
centralizer will be invariant under conjugation by members of $M;$
hence it is invariant under conjugation by
all members of $MK=G,$ i.e., it is normal in~$G.$
So by~\eqref{d.K_has_no_nm},
%
\begin{equation}\begin{minipage}[c]{35pc}\label{d.triv_cntrlzr}
The centralizer of $M$ in $K$ is trivial.
\end{minipage}\end{equation}

Returning to what we proved earlier about $H,$ note that
if $f^{-1}(K) < H$ had nontrivial intersection with
$H_0,$ then this would centralize $H_1,$ hence its image
in $G$ would be a subgroup of $K,$ nontrivial by~\eqref{d.triv_cap},
that centralized the image of $H_1,$ which we saw in~\eqref{d.M1,M2}
contains $M,$ contradicting~\eqref{d.triv_cntrlzr}.
So $f^{-1}(K)$ has trivial intersection with $H_0;$ and
similarly with $H_1.$
In view of~\eqref{d.H12}, this
forces $f^{-1}(K)$ to be the graph of an isomorphism between subgroups
$K_0 < G_0$ and $K_1 < G_1,$ each isomorphic to $K,$ and we see
that these act on $M_0$ and $M_1$ as $K$ acts on $M.$
Hence the
subgroup $M_0\,K_0 < G_0$ (and likewise $M_1\,K_1 < G_1)$ is isomorphic
to $MK = G,$ again contradicting our assumption that $G$ is not
a subquotient of $G_0$ or $G_1.$
This gives the desired contradiction to~(ii) for
$G$ not Rhodes-prime with respect to direct products.

To prove Rhodes-primeness with respect to direct products
in case~(iii), suppose a subgroup $H$ of a finite group $G_0\times G_1$
can be mapped surjectively to $G=\Z_{p^n}.$
Then an element mapping to a generator of $G$ must
have order divisible by $p^n.$
Since the order of an element of $G_0\times G_1$ is the least
common multiple of the orders of its components, one of those
components must have order divisible by $p^n,$
hence some power of that element will have order exactly $p^n,$
hence the subgroup of $G_0$ or $G_1$ that it generates will
be a subquotient of $G_0$ or~$G_1$ isomorphic to $\Z_{p^n}=G.$
\end{proof}

\end{document}
